%%%
%
% $Autor: Wings $
% $Date: 2024-10-31 11:15:45Z $
% $File: file.tex $
% $Version: 1 $
%
% Project: 
%
%%%


\chapter{\textcolor{red}{Hardware XY}}

\section{Jetson Nano}
\label{sec:jetsonnano}

\subsection{Description}
\label{subsec:description}

System-on-a-Chip (SoC) platforms provide a cost-effective and power-efficient solution for developing embedded machine learning systems. 
They are particularly suited for \textbf{inference tasks}, where pre-trained deep learning models are deployed to process real-world data in real time. 
Among such SoC systems, the \textbf{NVIDIA Jetson Nano} stands out as one of the most compact and capable AI modules available for edge computing. 
It belongs to NVIDIA’s Jetson family, which includes devices like Jetson Xavier NX and Jetson Orin Nano, but represents the most accessible entry point for developers building low-power AI vision systems \cite{jetsonNanoDatasheet, JetsonGettingStarted:2019}.

\bigskip

The Jetson Nano combines a \textbf{quad-core ARM Cortex-A57 CPU} with a \textbf{128-core Maxwell GPU}, offering approximately \textbf{0.5~TFLOPs (FP16)} of performance at a typical power draw of only 5–10~W. 
This makes it highly efficient for computer vision, robotics, and real-time deep learning inference. 
Its open development environment and support for the \textbf{JetPack SDK} (based on Ubuntu Linux) allow developers to integrate AI frameworks such as TensorFlow, PyTorch, OpenCV, and NVIDIA TensorRT easily.

\bigskip

From an industry standpoint, the Jetson Nano competes with other low-cost AI-capable platforms such as:
\begin{itemize}
	\item \textbf{Intel Neural Compute Stick~2:} a USB-based accelerator using the Movidius Myriad~X VPU, designed for embedded computer vision tasks and compatible with Raspberry~Pi and other Linux systems \cite{IntelStickDataSheet:2019}.
	\item \textbf{Google Coral Edge~TPU:} an AI accelerator board powered by Google's Edge TPU chip, offering optimized TensorFlow Lite inference for low-latency vision applications \cite{GoogleCoral:2019, Vincent:2018}.
	\item \textbf{Raspberry~Pi~5 + Hailo~AI Kit:} a more recent setup combining ARM CPUs with a small neural accelerator, but still limited in GPU throughput compared to the Jetson platform.
\end{itemize}

\bigskip

Other commercial SoCs, such as Apple’s A-series chips (A12~Bionic), Samsung’s Exynos~980 with integrated NPU, and Huawei’s Kirin~990, demonstrate similar edge AI capabilities but are proprietary and not intended for open embedded development \cite{SamsungExynos980:2020,HuaweiKirin990:2020}. 

\bigskip
The Jetson Nano, in contrast, remains \textbf{developer-oriented, open, and extensible}, offering a balance between computational performance, power efficiency, and cost (approximately 180–250~\texteuro{} in Germany). 
It serves as the foundation for the current project on \textbf{real-time license plate detection}, enabling efficient on-device image inference without the need for cloud processing. Overall, the Jetson Nano’s combination of GPU parallelism, open-source software ecosystem, and compatibility with AI frameworks makes it one of the most accessible and powerful tools for academic and industrial AI deployment at the edge.

%----------------------------------------------------------
\subsection{Applications of the Jetson Nano}
\label{subsec:applications}

The Jetson Nano can be utilized for a wide range of applications, ranging from basic computing tasks to complex edge-AI deployments. 
Its flexibility allows it to function both as a compact single-board computer and as a high-performance inference device for embedded vision systems \cite{Kurniawan:2021, JetsonProjects:2020}. 

\bigskip
Figure~\ref{fig:jetson_applications} provides a visual overview of some of the representative deep-learning and computer-vision tasks that can be executed efficiently on the Jetson Nano platform. 
These include \textit{image classification, object detection, semantic segmentation, pose estimation, action recognition, background removal,} and \textit{monocular depth estimation}. 
Such diverse capabilities demonstrate the Nano’s suitability for both research and industrial-level AI applications.

\begin{figure}[H]
	\centering
	\includegraphics[width=0.95\linewidth]{HardwareDesc/applications.jpeg}
	\caption{Representative AI and computer vision tasks deployable on the NVIDIA Jetson Nano platform \cite{JetsonProjects:2020}.}
	\label{fig:jetson_applications}
\end{figure}

\bigskip
\textbf{1. General Computing:}  
The Jetson Nano can perform conventional tasks such as running web servers, accessing cloud environments, developing software, or even low-end gaming. 
This versatility is due to its full Ubuntu~Linux environment and GPU support for graphical rendering.

\bigskip
\textbf{2. AI and Computer Vision:}  
The true potential of the Jetson Nano lies in its ability to deploy pre-trained neural networks for object detection, tracking, segmentation, and classification.  
The integrated \textbf{TensorRT SDK} optimizes AI models for fast inference, enabling applications such as:
\begin{itemize}
	\item Vehicle and pedestrian detection in smart cities.
	\item Facial recognition and identity verification.
	\item Edge-based anomaly detection in industrial systems.
\end{itemize}

\bigskip
\textbf{3. Robotics and Automation:}  
Community-driven projects, as showcased on NVIDIA’s Jetson developer portal, demonstrate the Nano’s use in autonomous vehicles, drones, and robotic arms.  
These projects leverage the Nano’s GPIO, I\textsuperscript{2}C, and CSI camera interfaces to build real-time perception and control loops \cite{JetsonProjects:2020}.

\bigskip
\textbf{4. Assistive Technologies:}  
In accessibility-focused research, the Jetson Nano has been used for projects that convert handwritten text to speech for visually impaired users or for real-time sign language recognition using CNN-based models trained in TensorFlow~2.0 on Google~Colab.

\bigskip
\textbf{5. Traffic Monitoring and Smart Surveillance:}  
Jetson Nano’s ability to process video streams in real time makes it ideal for traffic congestion analysis and license plate recognition systems — the focus of this project.  
Its edge-processing capability ensures privacy compliance and low latency compared to cloud-based analytics.

\bigskip
In summary, the Jetson Nano bridges the gap between experimental AI research and real-world embedded deployment. 
Its balance of performance, cost, and open-source accessibility enables both academic and industrial innovation in fields like automation, surveillance, and intelligent transportation systems.

\subsection{Construction}
\label{subsec:construction}

The complete hardware configuration is based on the \textbf{Seeed Studio reComputer J1010} system, which integrates the \textbf{NVIDIA Jetson Nano (4 GB, B01 revision)} module on the \textbf{J101 carrier board}. 
The Jetson Nano module itself measures approximately 69.6 mm $\times$ 45 mm and is dominated by an aluminum heatsink that covers the NVIDIA Tegra X1 system-on-chip.  
Hidden beneath the heatsink, the Tegra X1 combines a quad-core ARM Cortex-A57 CPU cluster and a 128-core Maxwell GPU fabricated in a 20 nm process, providing up to 0.5 TFLOPs (FP16) of compute performance at a typical power draw of 5–10 W \cite{jetsonNanoDatasheet}.  
The module connects to the carrier board through a 260-pin SODIMM edge connector, exposing all high-speed I/O lines.

\bigskip
\textbf{Physical Layout:}  
The carrier board has approximate dimensions of 100 mm $\times$ 80 mm $\times$ 29 mm and is mounted inside an anodized aluminum case supplied by Seeed Studio.  
Figure \ref{fig:recomputer_layout} (official top-view diagram) labels the major connectors, including the USB-C power port (J1), HDMI output (J3), dual USB 2.0 Type-A ports (J5), USB 3.0 Type-A (J4), RJ-45 Ethernet (J6), MIPI-CSI camera connectors (J9/J13), and 40-pin GPIO header (J10).  
A passive heatsink ensures adequate cooling for continuous inference workloads; an optional 5 V PWM fan can be attached to the J8 header for active thermal control.

\bigskip
\textbf{Core Components:}
\begin{itemize}
	\item \textbf{Compute Module:} NVIDIA Jetson Nano (B01 SoC, 4 GB LPDDR4 memory, 16 GB eMMC 5.1 flash)
	\item \textbf{Carrier Board:} Seeed Studio J101 with USB, HDMI, CSI, M.2 Key E, Ethernet, and 40-pin GPIO
	\item \textbf{Power Supply:} 5 V DC $\pm$ 0.25 V, 4 A via USB-C input (J1); reverse-current protection integrated
	\item \textbf{Storage:} microSD card (64–128 GB, UHS-I recommended) used for OS image and dataset storage
	\item \textbf{Networking:} RJ-45 Gigabit Ethernet port (10/100/1000 BASE-T)
	\item \textbf{Camera:} MIPI CSI-2 camera (IMX219, 8 MP, 4-lane)
	\item \textbf{Cooling:} Passive aluminum heatsink; optional PWM-controlled 5 V fan (J8)
\end{itemize}

\begin{figure}[H]
	\centering
	\includegraphics[width=0.9\linewidth]{HardwareDesc/reComputerJ101v2.png}
	\caption{Top view of the Seeed Studio reComputer J1010 showcasing major connectors and components, including power input, USB ports, HDMI, CSI camera slots, and the 40-pin GPIO header \cite{seeedJ101Manual}.}
	\label{fig:recomputerJ101v2}
\end{figure}

\bigskip
\textbf{Mechanical Arrangement:}  
The Jetson Nano module is fixed to the carrier through four M2.5 mounting screws.  
The board assembly is then secured inside the aluminum enclosure with standoffs that provide isolation and airflow channels.  
Cut-outs on the side panels align precisely with the external I/O ports, allowing straightforward access during prototyping and debugging.

\bigskip
\textbf{Power Architecture:}  
The J101 board features a single 5 V rail that powers the SoC, USB hub, and Ethernet PHY through onboard DC-DC converters.  
An internal power-management IC monitors undervoltage events and initiates a safe power-down sequence.  
Average consumption during neural-network inference (YOLOv8 Tiny model, 30 FPS 1080p input) is approximately 7.5 W, measured at 5 V input.  
Peak loads of 9–10 W can occur during model initialization, which justifies the use of a 4 A supply.

\bigskip
\textbf{Peripheral Integration:}  
The reComputer’s J10 40-pin header maintains the Raspberry Pi-compatible pinout, permitting connection of sensors, relays, and GPIO-controlled modules through standard ribbon cables.  
Two CSI camera ports (J9 and J13) enable stereo vision or dual-camera processing, though only one port is utilized in this project for license-plate detection.  
The M.2 Key E slot supports Wi-Fi and Bluetooth modules through PCIe x1 and USB 2.0 lanes.  
The RTC header (J18) can house a CR2032 cell to preserve system time when powered off.  

\bigskip
\textbf{Connectivity and Display:}  
For configuration and monitoring, a display may be attached to the HDMI 2.0 port (J3) supporting resolutions up to 3840 × 2160 @ 60 Hz.  
During unattended operation, the system runs headless and communicates over Ethernet SSH.  
Gigabit Ethernet provides sufficient bandwidth for remote dataset transfer and model updates.  

\bigskip
\textbf{Storage and Operating System:}  
The microSD card serves as the main boot device and storage medium.  
JetPack 5.x (derived from Ubuntu 20.04 LTS) is flashed onto the card using NVIDIA’s SDK Manager.  
This OS package includes CUDA 10.2, cuDNN, TensorRT, and OpenCV 4.x libraries pre-optimized for the Tegra X1 architecture.

\bigskip
\textbf{Thermal and Environmental Considerations:}  
Under sustained GPU utilization, SoC temperature typically stabilizes between 55 °C and 65 °C with the passive heatsink.  
In warmer environments (> 30 °C ambient), a 5 V PWM fan connected to J8 is recommended.  
The board operates reliably between 0 °C and 40 °C ambient and at relative humidity up to 90 \% non-condensing.

\bigskip
\textbf{System Assembly in the Project:}  
For the License Plate Detection application, the Jetson Nano is mounted on an acrylic base plate with vibration-absorbing spacers.  
The CSI camera is fixed on an adjustable stand at a 20–25° angle relative to vehicle direction, ensuring clear plate capture at a 1.5 m height.  
Ethernet provides Internet access for model updates and data synchronization with the host laptop used for training.

\bigskip
Overall, the modular construction of the Seeed Studio reComputer J1010 simplifies expansion, maintenance, and cooling, making it a reliable platform for continuous AI inference at the edge \cite{seeedJ101Manual,edgeAIbook}.


\subsection{Interfaces}
The Seeed Studio J101 carrier board exposes multiple connectors for communication, expansion, and debugging. 
Each interface corresponds to the labeled connector on the PCB as defined in the official datasheet \cite{seeedJ101Manual}.

\subsubsection{J10 – 40-Pin Expansion Header (GPIO)}
Provides 3.3~V general-purpose I/O pins supporting digital I/O, I\textsuperscript{2}C, SPI, UART, I\textsuperscript{2}S, and PWM. 
Compatible with the Jetson.GPIO Python library for user-space control.

\subsubsection{J14 – Serial and Button Header (UART / Debug)}
Exposes UART2 (TX/RX) for serial debugging, along with SYS\_RESET, PWR\_BTN, and FORCE\_RECOVERY pins for power and firmware control. 
Auto power-on can be disabled by shorting pins~5--6.

\subsubsection{J1–J5 – USB Interfaces}
\begin{itemize}
	\item \textbf{J1:} USB-C power input (5~V, up to 4~A)
	\item \textbf{J2:} USB-C data/recovery port
	\item \textbf{J4:} USB~3.0 Type-A (5~Gbps)
	\item \textbf{J5:} Dual USB~2.0 Type-A (480~Mbps total, 2~A shared VBUS)
\end{itemize}

\subsubsection{J6 – Gigabit Ethernet Interface}
Standard RJ-45 connector supporting 10/100/1000~BASE-T with link (green) and activity (yellow) LEDs. 
Optional Power-over-Ethernet (PoE) is supported via additional circuitry.

\subsubsection{J3 – HDMI Interface}
HDMI~2.0 Type-A connector supporting up to 4K~@~60~Hz video output with audio, CEC, and DDC communication.

\subsubsection{J11 – M.2 Key E Socket}
Provides PCIe~x1, USB~2.0, UART, I\textsuperscript{2}S, and I\textsuperscript{2}C interfaces for Wi-Fi/Bluetooth modules.

\subsubsection{J9 \& J13 – Camera (CSI) Interfaces}
Two 15-pin MIPI CSI-2 connectors supporting dual 4-lane cameras. 
Controlled through I\textsuperscript{2}C lines for power and sensor configuration.

\subsubsection{J8 – Fan Connector}
4-pin PWM header (5~V, PWM, Tach, GND) for optional active cooling. 
PWM control is available via GPIO14, and tachometer feedback via GPIO08.

\subsubsection{J18 – RTC Coin Cell Battery Holder}
Provides backup power for the Real-Time Clock (RTC) using a CR2032 cell.

\subsubsection{J19 – SDIO Interface}
Dedicated microSD interface supporting 1.8~V/3.3~V operation for SD cards or SDIO peripherals.

\subsubsection{Summary of Jetson Nano Interfaces}
\begin{table}[H]
	\centering
	\caption{Summary of Jetson Nano Interfaces (reComputer J1010)}
	\begin{tabular}{|p{1.5cm}|p{3cm}|p{4cm}|p{2.5cm}|p{3cm}|}
		\hline
		\textbf{Header / Port} & \textbf{Type / Connector} & \textbf{Key Functions} & \textbf{Voltage / Speed} & \textbf{Remarks} \\
		\hline
		J1 & USB-C & Power Input & 5~V~$\pm$~0.25~V, 4~A & Main power supply \\
		J2 & USB-C & Data / Recovery & USB~2.0 & Device mode \\
		J3 & HDMI Type-A & Display Output & HDMI~2.0 (4K@60~Hz) & With audio \\
		J4 & USB~3.0 Type-A & Host Port & 5~Gbps & Camera or SSD \\
		J5 & Dual USB~2.0 Type-A & Peripheral Ports & 480~Mbps & Keyboard, Mouse \\
		J6 & RJ-45 Ethernet & Networking & 1~Gbps & PoE optional \\
		J8 & 4-Pin Fan & Cooling Control & 5~V PWM & Optional fan \\
		J9 / J13 & 15-Pin CSI & Camera Input & 4-Lane MIPI & Dual camera support \\
		J10 & 40-Pin Header & GPIO / I$^2$C / SPI / UART / I$^2$S & 3.3~V logic & Expansion header \\
		J11 & M.2 Key E & PCIe~x1, USB2, I$^2$C & 3.3~V & Wireless modules \\
		J14 & Button / UART & Power \& Debug Control & 3.3~V & System interface \\
		J18 & Coin Cell Holder & RTC Backup & 3~V & CR2032 battery \\
		J19 & SDIO Interface & Storage Interface & 1.8~V/3.3~V & microSD slot \\
		\hline
	\end{tabular}
\end{table}

\subsubsection{ External Interfaces of Jetson Nano}

\begin{figure}[H]
	\centering
	\includegraphics[width=0.9\linewidth]{HardwareDesc/jetson-side.png}
	\caption{Front I/O panel of the Seeed Studio reComputer J1010 with labeled external interfaces (1--6) corresponding to Table~\ref{tab:j1010_interfaces} \cite{seeedJ101Manual}.}
	\label{fig:jetson_side}
\end{figure}


\begin{table}[H]
	\centering
	\caption{External interfaces of the Seeed Studio reComputer J1010}
	\label{tab:j1010_interfaces}
	\begin{tabular}{|c|l|l|}
		\hline
		\textbf{No.} & \textbf{Name} & \textbf{Note} \\ \hline
		1 & DCIN (Type-C interface) & Power supply only \\ \hline
		2 & HDMI & \\ \hline
		3 & 1x USB 3.0 Type-A port & \\ \hline
		4 & 2x USB 2.0 Type-A ports & \\ \hline
		5 & LAN & \\ \hline
		6 & USB Type-C port & Data transmitted only \\ \hline
	\end{tabular}
\end{table}


\subsection{Pin Configuration}
All general-purpose I/O pins on J10 operate at \textbf{3.3~V logic levels}. 
They support configurable modes such as digital I/O, PWM, SPI, UART, I\textsuperscript{2}C, and I\textsuperscript{2}S. 
Over-voltage above 3.3~V can damage the SoC pins. 
Pin control and configuration are performed through the Jetson.GPIO library or the device tree \cite{seeedJ101Manual}.

\subsection{Data Quality}
The Jetson Nano offers consistent high-resolution image capture through its dedicated MIPI CSI-2 interface, minimizing latency compared to USB cameras. 
GPU-accelerated image preprocessing ensures robust frame quality and stability for license plate detection under variable illumination \cite{edgeAIbook}.

\subsection{Data Quantity}
The CSI camera streams up to 30~FPS at 720p or 1080p resolution, producing 2--4~MB per frame. 
A 64--128~GB microSD card or USB~3.0 SSD provides sufficient capacity for long-duration data recording.

\subsection{Specification of Jetson Nano}
\begin{table}[H]
	\centering
	\caption{Technical Specifications of NVIDIA Jetson Nano (B01 Revision)}
	\begin{tabular}{|p{5cm}|p{8cm}|}
		\hline
		\textbf{Parameter} & \textbf{Specification} \\
		\hline
		GPU & 128-core NVIDIA Maxwell architecture \\
		CPU & Quad-core ARM Cortex-A57 MPCore \\
		Memory & 4~GB LPDDR4 @ 1600~MHz (25.6~GB/s) \\
		Storage & 16~GB eMMC + microSD expansion \\
		Camera Support & Dual CSI-2 (4-lane each) \\
		Networking & Gigabit Ethernet (10/100/1000~Mbps) \\
		Video Decode & Up to 4K~@~60~Hz (HEVC/H.264) \\
		Video Encode & Up to 4K~@~30~Hz \\
		Power Input & 5~V DC, up to 4~A (USB-C) \\
		Typical Power Use & 5--10~W \\
		Operating System & NVIDIA JetPack 5.x (Ubuntu 20.04 LTS) \\
		\hline
	\end{tabular}
\end{table}

\subsection{Constraints}
\begin{itemize}
	\item Limited 4~GB RAM restricts execution of multiple large AI models concurrently.
	\item All GPIO pins operate strictly at 3.3~V logic; higher voltages are not supported.
	\item No native NVMe slot; high-speed storage achievable via USB~3.0 SSD.
	\item Continuous GPU inference may require active cooling.
	\item Stable 5~V~$\pm$~5\% power supply is required to prevent brown-out resets.
\end{itemize}

\subsection{Dimension of Jetson Nano}
\begin{itemize}
	\item Module Size: 69.6~mm $\times$ 45~mm
	\item Carrier Board Size: 100~mm $\times$ 80~mm
	\item Enclosure Size: 130~mm $\times$ 120~mm $\times$ 50~mm
	\item Total Weight: approximately 300~g (including case and heatsink)
\end{itemize}

\subsection{Simple Code for Jetson Nano Reference Carrier Board}
The following Python program demonstrates GPIO functionality by blinking an LED connected to pin~12 on the J10 header.

\begin{lstlisting}[language=Python, caption={LED Blink Test on Jetson Nano}]
	import Jetson.GPIO as GPIO
	import time
	
	LED_PIN = 12
	GPIO.setmode(GPIO.BOARD)
	GPIO.setup(LED_PIN, GPIO.OUT, initial=GPIO.LOW)
	
	try:
	while True:
	GPIO.output(LED_PIN, GPIO.HIGH)
	time.sleep(0.5)
	GPIO.output(LED_PIN, GPIO.LOW)
	time.sleep(0.5)
	except KeyboardInterrupt:
	GPIO.cleanup()
\end{lstlisting}

This code verifies GPIO operation and confirms the Jetson Nano’s correct hardware configuration under JetPack~5.x.